\chapter{Valuación}\label{ch:valuation}

% Alcance: DCF, múltiplos (EV/EBITDA, P/E, EV/Ingresos), valor terminal, ponderación por escenarios.
% Referencia cruzada: la fórmula del valor terminal es una aplicación de la perpetuidad creciente --
%   \cref{eq:gordon} en \cref{ch:corporate-finance-core}.

\section{Flujo de caja descontado}\label{sec:dcf}
% complejidad: 5 -> figura dcf-bridge

¿Cuánto vale hoy la empresa en su totalidad? El DCF responde convirtiendo cada flujo de caja libre futuro
--- a lo largo de un período de pronóstico explícito y de un valor terminal más allá de él --- en un único
valor presente al costo de capital. La respuesta determina si conviene invertir, adquirir o vender. Una
empresa es un conjunto de derechos futuros sobre efectivo: el FCF de cada año se descuenta individualmente,
exactamente como en \cref{eq:pv}; el valor terminal captura todos los flujos de caja más allá del horizonte
de pronóstico como una suma única. En la práctica, el valor terminal casi siempre domina, razón por la cual
los supuestos de crecimiento que lo sustentan merecen mayor escrutinio que los pronósticos del período
explícito.

El flujo de caja libre es el efectivo generado después de impuestos (calculados sobre el EBIT, lo que
preserva el escudo de depreciación) y de cubrir las necesidades de reinversión:
\[
  \mathrm{FCF}_t \;=\; \mathrm{EBITDA}_t \;-\; \text{taxes}_t
                       \;-\; \Delta\mathrm{WC}_t \;-\; \mathrm{CapEx}_t.
\]
Se descuenta el FCF de cada año al WACC (\cref{eq:wacc}) y se suma el valor presente del valor terminal
\(\mathrm{TV}\) en el año horizonte \(n\):
\begin{equation}\label{eq:dcf}
  EV \;=\; \sum_{t=1}^{n}\frac{\mathrm{FCF}_t}{(1+\mathrm{WACC})^{t}}
           \;+\; \frac{\mathrm{TV}}{(1+\mathrm{WACC})^{n}}.
\end{equation}
El WACC (\cref{eq:wacc}) se deriva en \cref{ch:risk-and-capital-structure}; su componente de \emph{equity}
es el costo de \emph{equity} del CAPM (\cref{eq:capm}). La fórmula del valor terminal se desarrolla en
\cref{sec:terminal-value}. Nótese que \cref{eq:dcf} es la extensión multiperiodo de \cref{eq:npv} en
\cref{ch:corporate-finance-core}: el desembolso inicial \(I_0\) es el precio de adquisición, y el valor
empresarial es la suma que hace \(\mathrm{NPV}=0\).

El DCF vale únicamente lo que valen sus insumos. El WACC se estima, no se observa; los pronósticos de FCF
son narrativas disfrazadas de números; y el valor terminal --- típicamente
\qty{60}{\percent}--\qty{80}{\percent} del EV --- hereda todos los supuestos del período explícito y los
extrapola indefinidamente. La sensibilidad al WACC es no lineal: un aumento de un punto porcentual en la
tasa de descuento puede reducir el valor empresarial en \qty{20}{\percent}--\qty{30}{\percent} en empresas
de alto crecimiento, donde la mayor parte del valor reside en el período terminal.

\begin{example}[DCF a cinco años para el B2B SaaS de referencia]
Cinco años explícitos de pronóstico para el B2B SaaS de referencia. El FCF del año 1 es \money{600000},
con crecimiento del \qty{30}{\percent} en los años 1--3 (fase de ajuste producto-mercado), que se
desacelera al \qty{10}{\percent} en los años 4--5. El WACC es \qty{11}{\percent} (\cref{eq:wacc}
en \cref{ch:risk-and-capital-structure}; equivale al costo de \emph{equity} porque el negocio
se financia íntegramente con capital propio).

\[
\begin{array}{lrrrrrr}
  & \mathrm{Y1} & \mathrm{Y2} & \mathrm{Y3} & \mathrm{Y4} & \mathrm{Y5} \\[2pt]
  \mathrm{FCF} & \num{600000} & \num{780000} & \num{1014000} & \num{1115400} & \num{1226940} \\
  \text{Factor de descuento} & 0.901 & 0.812 & 0.731 & 0.659 & 0.593 \\
  \mathrm{PV} & \num{540541} & \num{633246} & \num{741379} & \num{735047} & \num{727612} \\
\end{array}
\]

Suma de los valores presentes del período explícito: \(\approx\money{3377825}\).
Valor terminal (FCF normalizado del año 5 con tasa de crecimiento estable del \qty{3}{\percent};
calculado íntegramente en \cref{sec:terminal-value}): \(\approx\money{9400000}\) VP.
Valor empresarial: \(\money{3377825}+\money{9400000}\approx\money{12800000}\).

Implicación gerencial: los años explícitos aportan solo el \qty{26}{\percent} del EV total.
Elevar el supuesto de crecimiento terminal del \qty{3}{\percent} al \qty{4}{\percent} añade más
de \money{1000000}; reducir el WACC un punto agrega aún más. Los inversores que se concentran
en el ingreso de corto plazo pasan por alto dónde reside realmente el valor.
\end{example}

\begin{figure}[htbp]
  \centering
  \includegraphics[width=0.86\linewidth]{dcf-bridge}
  \caption{Cascada DCF: valor presente de cada FCF del período explícito apilado frente al valor presente del valor
  terminal. El valor terminal domina (\(\approx\)\qty{74}{\percent} del valor empresarial),
  lo que ilustra por qué los supuestos de tasa de crecimiento y tasa de descuento en el período
  terminal superan en importancia a la precisión del pronóstico de corto plazo.}
  \label{fig:dcf-bridge}
\end{figure}

\begin{admonition}{Advertencia}
Tratar el resultado del DCF como \emph{el} valor en lugar de uno de los escenarios posibles. El modelo
responde «¿cuál es el valor \emph{dados estos supuestos}?» --- no «¿cuánto pagará un comprador?»
Un DCF preciso construido sobre crecimiento optimista y un WACC bajo es una forma sofisticada de confirmar
una convicción preexistente. La disciplina consiste en ejecutar el modelo con supuestos que se defenderían
ante un comité de crédito, no con los supuestos que hacen funcionar la operación.
\end{admonition}

El DCF apalancado descuenta los flujos de caja libres del \emph{equity} directamente al costo de \emph{equity};
el DCF no apalancado descuenta los FCF de la firma al costo no apalancado y luego deduce la deuda. El valor
presente ajustado (APV) separa el valor de la firma no apalancada del valor presente del escudo fiscal, lo cual
resulta útil cuando el apalancamiento cambia materialmente a lo largo del horizonte de pronóstico --- como en
una adquisición apalancada (LBO) o una construcción de infraestructura por etapas.

El DCF formaliza la misma lógica que la regla del NPV en \cref{ch:corporate-finance-core}
(\cref{eq:npv}) y la fórmula del CLV en \cref{ch:customer-economics} (\cref{eq:clv}):
todos son valores presentes de flujos de caja futuros descontados al costo de capital apropiado.
\textcite{brealey2020}, \textcite{damodaran2012} y
\textcite{berk6thcorporate} desarrollan cada uno el \emph{framework} del DCF; \textcite{damodaran2012}
es el tratamiento más riguroso de la mecánica del valor terminal.

\section{Valor terminal}\label{sec:terminal-value}
% complejidad: 4

Más allá del horizonte de pronóstico explícito, la empresa no deja de existir. El valor terminal convierte
los flujos de caja posteriores al horizonte en un único número anclado en el año \(n\), sin necesidad de
extender una hoja de cálculo de veinte años. Al llegar al horizonte se asume que la empresa alcanza un estado
estable: el crecimiento se estabiliza en una tasa de largo plazo \(g\) y los márgenes dejan de expandirse. La
fórmula de la perpetuidad creciente (\cref{eq:gordon}) valora ese estado estable; descontarla al presente
arroja la contribución del valor terminal al valor empresarial.

Sea \(\mathrm{FCF}_{n+1}\) el primer flujo de caja posterior al horizonte (el FCF del año \(n\) crecido un
período) y \(g\) la tasa de crecimiento de la perpetuidad. El valor terminal en el año \(n\) es una aplicación
directa de \cref{eq:gordon}:
\begin{equation}\label{eq:terminal-value}
  \mathrm{TV} \;=\; \frac{\mathrm{FCF}_{n+1}}{\mathrm{WACC}-g}.
\end{equation}
Por eso \cref{eq:terminal-value} está implícita en todo DCF: al sustituir
\(\mathrm{TV}/(1+\mathrm{WACC})^n\) en \cref{eq:dcf}, la fórmula completa se colapsa en una perpetuidad
creciente de dos etapas. La fortaleza de marca eleva \(g\) al horizonte
(\cref{ch:brand-growth}), que se multiplica a través de \cref{eq:terminal-value} en el valor empresarial.

El valor terminal es mecánicamente correcto pero económicamente frágil. \(g\) debe ser inferior al WACC
(condición de convergencia de \cref{eq:gordon}); también debe ser inferior al crecimiento nominal del PIB
a largo plazo, ya que ninguna empresa puede crecer más que toda la economía de forma indefinida. Fijar
\(g=\qty{5}{\percent}\) con un WACC de \qty{11}{\percent} es defendible; fijar
\(g=\qty{9}{\percent}\) hace que el denominador se aproxime a cero y el valor terminal explote.

\begin{example}[Valor terminal para el SaaS de referencia]
El FCF normalizado del año 5 del SaaS es \money{1226940}; aplicando un período de crecimiento
estable (\qty{3}{\percent}) se obtiene \(\mathrm{FCF}_{6}=\money{1263748}\).
Con WACC de \qty{11}{\percent}:
\[
  \mathrm{TV} \;=\; \frac{\money{1263748}}{0.11-0.03} \;=\; \money{15796850}.
\]
Descontando ese valor terminal del año 5 cinco años atrás al \qty{11}{\percent}, se obtiene su
contribución presente: \(\mathrm{PV}(\mathrm{TV})=\money{15796850}/1.11^5\approx\money{9400000}\).
Sumado a los \money{3377825} de los valores presentes del período explícito, el valor empresarial es
\(\approx\money{12800000}\); el valor terminal representa el \qty{74}{\percent} de ese total ---
lo que confirma que la valuación es en gran medida una apuesta al estado estable, no a los
próximos cinco años.
\end{example}

\begin{admonition}{Advertencia}
\emph{El doble conteo del crecimiento} es el error canónico: incorporar crecimiento elevado en el período
explícito \emph{y} en la perpetuidad terminal. La \(g\) terminal debe reflejar únicamente el período
\emph{posterior} a los años explícitos, una vez que la dinámica competitiva ha comprimido los márgenes
hacia el estado estable del sector. Reutilizar la misma tasa de crecimiento del \qty{30}{\percent}
del año 1 en el valor terminal produce un absurdo. La misma trampa se señala para \cref{eq:gordon}
en \cref{sec:growing-perpetuity}.
\end{admonition}

Una metodología alternativa para el valor terminal es el enfoque de múltiplo de salida: estimar el
múltiplo de EBITDA que pagaría un comprador al horizonte y multiplicarlo por el EBITDA del año \(n\).
Es importante señalar que este enfoque no es independiente del DCF --- un múltiplo es en sí mismo una
perpetuidad creciente comprimida (\cref{sec:multiples}), por lo que ambos métodos convergen cuando sus
supuestos de crecimiento y descuento incorporados son consistentes. Utilizarlos como verificación
cruzada, y no como valuaciones independientes, es la mejor práctica.

El valor terminal es el componente dominante del valor empresarial de cualquier empresa cuyo valor
principal reside más allá del horizonte de pronóstico explícito --- es decir, la mayoría de las empresas
en etapa de crecimiento. \textcite{damodaran2012} ofrece el tratamiento más sistemático; \textcite{brealey2020}
cubre las condiciones de consistencia entre los enfoques de múltiplo de salida y perpetuidad.

\section{Múltiplos}\label{sec:multiples}
% complejidad: 3

Se busca una valuación rápida anclada en el mercado sin construir un DCF completo. Los múltiplos ofrecen
un atajo --- pero solo si se comprenden qué supuestos de crecimiento y rentabilidad están comprimiendo.
Malinterpretar esa compresión lleva a pagar el precio equivocado. Un múltiplo EV/Ingresos parece una
opinión del mercado, pero es un DCF disfrazado: el múltiplo incorpora supuestos sobre márgenes, crecimiento
y tasas de descuento. Entender el álgebra permite auditar si un múltiplo negociado está barato o caro en
relación con los fundamentos, y explicar por qué los múltiplos se contrajeron de 2021 a 2023.

Se parte de \cref{eq:gordon}: \(EV = \mathrm{FCF}/(r-g)\). Se expresa el FCF como margen sobre
los ingresos \(R\): \(\mathrm{FCF}=m\cdot R\). Se dividen ambos lados entre \(R\) y se permite un período
de crecimiento para que \(\mathrm{FCF}_1 = m\cdot R\cdot(1+g)\):
\begin{equation}\label{eq:ev-revenue}
  \frac{EV}{R} \;=\; \frac{m\,(1+g)}{\mathrm{WACC}-g}.
\end{equation}
Este es el múltiplo de ingresos consistente con el DCF. Aumenta con el margen \(m\), aumenta con el
crecimiento \(g\), y \emph{cae con el WACC} --- razón por la cual el alza de las tasas de interés
comprime los múltiplos de forma mecánica, independientemente de cualquier cambio en el negocio subyacente.
\cref{eq:ev-revenue} asume márgenes y crecimiento en estado estable; los múltiplos negociados también
incorporan sentimiento, primas de liquidez y comportamiento de rebaño entre pares. Un múltiplo sin una
posición sobre \(m\), \(g\) y el WACC no es una valuación --- es un precio, que puede o no reflejar el valor.

\begin{example}[Múltiplo de ingresos versus DCF para el SaaS de referencia]
El SaaS tiene un margen FCF del \qty{20}{\percent}, \(g=\qty{3}{\percent}\), WACC \qty{11}{\percent}.
De \cref{eq:ev-revenue}:
\[
  \frac{EV}{R} \;=\; \frac{0.20\times1.03}{0.11-0.03} \;=\; \frac{0.206}{0.08} \;=\; 2.6\times.
\]
Con ARR de \money{4000000}:
\[
  EV \;=\; 2.6\times\money{4000000} \;=\; \money{10400000}.
\]
El DCF de \cref{sec:dcf} arroja \(\approx\money{12800000}\) --- la
diferencia de \money{2400000} es la \emph{prima de crecimiento} que los múltiplos en estado estable
no capturan: el modelo usa solo \(g=\qty{3}{\percent}\) (largo plazo), pero la empresa está creciendo
actualmente al \qty{30}{\percent}. Un múltiplo calibrado con pares que crecen al
\qty{30}{\percent} daría un resultado significativamente mayor. La disciplina consiste en entender
qué tasa de crecimiento está incorporada antes de citar un múltiplo como evidencia de valor.
\end{example}

En 2021, el SaaS público cotizaba a 20--30$\times$ los ingresos; a finales de 2023 los múltiplos se habían
comprimido a 5--8$\times$. El motor dominante no fue el deterioro de la economía unitaria, sino una
repriorización mecánica: la tasa libre de riesgo subió del \qty{0}{\percent} al \qty{5}{\percent},
elevando el WACC y colapsando el denominador de \cref{eq:ev-revenue}. Si los múltiplos de 2021
reflejaban euforia irracional o si los de 2023 corrigieron en exceso es un asunto debatido; lo que no
está en disputa es que \cref{eq:ev-revenue} rastrea la dirección y la magnitud aproximada de la
compresión (\cref{ch:risk-and-capital-structure}).

El múltiplo EV/EBITDA elimina la estructura de capital y la política de depreciación, haciendo la
comparación entre empresas más limpia que el P/E. Su análogo en el DCF reemplaza \(m\cdot R\) con
EBITDA\(\times(1-t)\times(1-\mathrm{reinvestment\,rate})\) --- económicamente idéntico una vez que
los impuestos y la reinversión son consistentes. La ventaja práctica de los múltiplos es la rapidez y
el anclaje en el mercado; el riesgo es confundir un precio de mercado con un modelo de valor intrínseco.

Los múltiplos se tratan mejor como una verificación de cordura sobre un DCF, no como un reemplazo del mismo.
\textcite{damodaran2012} descompone sistemáticamente el álgebra detrás de cada múltiplo habitual;
el paralelismo con el modelo de Gordon (\cref{eq:gordon}) se hace explícito
en su análisis del ratio PEG y la inversión del rendimiento por dividendo.

\section{Análisis de escenarios}\label{sec:scenario-analysis}
% complejidad: 3

Un DCF de estimación puntual o un múltiplo implica certidumbre sobre un futuro incierto. El análisis de
escenarios reemplaza esa ilusión con una distribución de valores empresariales ponderada por probabilidades,
haciendo el riesgo explícito y la decisión defendible. En lugar de preguntar «¿cuál es el valor?», el
análisis de escenarios pregunta «¿cuál es el valor \emph{condicionado a cada futuro plausible}?» Se ponderan
los resultados por sus probabilidades y el valor esperado captura tanto el potencial alcista como el
bajista en un único número procesable.

Sea \(p_i\) la probabilidad del escenario \(i\) y \(EV_i\) el valor empresarial bajo ese escenario
(calculado mediante \cref{eq:dcf} o un múltiplo de \cref{eq:ev-revenue}). El EV ponderado por escenarios es:
\begin{equation}\label{eq:scenario-ev}
  EV^* \;=\; \sum_{i} p_i \cdot EV_i.
\end{equation}
\cref{eq:scenario-ev} reutiliza la maquinaria del valor esperado de \cref{eq:ev} en
\cref{ch:decisions-under-uncertainty}: es la misma ponderación por probabilidades aplicada a valuaciones
en lugar de pagos. En el \emph{framework} de ese capítulo, un tomador de decisiones con aversión al riesgo
ajustaría adicionalmente a la baja por la varianza; aquí el ajuste está incorporado en la tasa de descuento
en lugar de en las probabilidades. El resultado depende enteramente de la calidad de las estimaciones de
\(p_i\) y \(EV_i\). Las probabilidades de los escenarios rara vez son objetivas --- reflejan las
convicciones previas del analista. La fórmula disciplina la \emph{estructura} del pensamiento (tres futuros,
asignar pesos, calcular) pero no puede sustituir al buen juicio sobre qué futuros son plausibles.

\begin{example}[Decisión de adquisición ponderada por escenarios]
Tres escenarios para el SaaS ante una posible decisión de adquisición:

\begin{description}
  \item[Bajista (\(p=0.25\))] La competencia se intensifica; el crecimiento se estanca; EV \money{8000000}.
  \item[Base (\(p=0.50\))] La trayectoria actual se mantiene; EV \money{15000000}.
  \item[Alcista (\(p=0.25\))] Los efectos de red de la plataforma se activan; EV \money{28000000}.
\end{description}

Aplicando \cref{eq:scenario-ev}:
\[
  \begin{aligned}
    EV^* &= 0.25\times\money{8000000}
          + 0.50\times\money{15000000}
          + 0.25\times\money{28000000} \\
         &= \money{2000000}+\money{7500000}+\money{7000000}
          = \money{16500000}.
  \end{aligned}
\]
El EV esperado es \money{16500000} --- superior al DCF de estimación puntual de
\money{12800000} porque el escenario alcista se pondera igual que el bajista pero se sitúa
más lejos de la base.
La asimetría entre escenarios determina el resultado; la sensibilidad a las asignaciones de \(p_i\)
importa más que la precisión del modelo dentro de cada escenario.
\end{example}

\begin{admonition}{Advertencia}
Colapsar tres escenarios en un único «caso base» para evitar el malestar del escenario bajista.
\emph{La ponderación por probabilidades no es pesimismo} --- es contabilizar el rango completo de
resultados. Un escenario bajista con \(p=0.25\) no es un peor caso que se puede descartar;
es la afirmación de que una de cada cuatro decisiones similares termina ahí.
\end{admonition}

La simulación Monte Carlo extiende el análisis de escenarios extrayendo de distribuciones de probabilidad
continuas sobre cada insumo (WACC, \(g\), FCF del año 1) y calculando la distribución completa del EV.
El resultado es un histograma de valores empresariales en lugar de tres estimaciones puntuales, lo que
revela colas gruesas que los pesos de los escenarios pueden ocultar. La pregunta estructural --- qué
insumos generan la mayor varianza --- se responde mejor mediante un análisis de sensibilidad de Sobol,
que descompone la varianza total del EV atribuible a cada insumo.

El análisis de escenarios enlaza con el \emph{framework} de decisiones de \cref{ch:decisions-under-uncertainty}
y avanza hacia la elección de estructura de capital de \cref{ch:risk-and-capital-structure}:
una mayor varianza en el EV argumenta a favor de un menor apalancamiento, dado que los costos de
dificultades financieras resultan más probables. \textcite{damodaran2012} desarrolla la extensión de
simulación; \textcite{brealey2020} cubre el enfoque de análisis de escenarios de punto de equilibrio
como complemento.
